\documentclass[11pt]{article}

\usepackage{amsmath,amssymb,amsthm,mathtools}
\usepackage[margin=1in]{geometry}
\usepackage{microtype}
\usepackage{hyperref}
\usepackage{enumitem}
\usepackage{booktabs}
\hypersetup{colorlinks=true, linkcolor=blue, citecolor=blue, urlcolor=blue}

\newtheorem{theorem}{Theorem}
\newtheorem{proposition}{Proposition}
\newtheorem{corollary}{Corollary}
\newtheorem{lemma}{Lemma}
\newtheorem{remark}{Remark}
\newtheorem{definition}{Definition}

\newcommand{\Zm}{\mathbb Z_m}
\newcommand{\I}{\mathbf 1}
\newcommand{\sgn}{\operatorname{sgn}}

\title{Toward the even-m D5 case\\
Parity barrier, formal extension theorem, and a critical-row program}
\author{}
\date{14 March 2026}

\begin{document}
\maketitle

\begin{abstract}
This note records what can already be proved about the even-modulus D5 case
from the stabilized odd-modulus boundary theorem and the d=3 even-modulus
philosophy. Two points are rigorous. First, the Kempe sign-product argument from
$d=3$ extends verbatim to $D_5(m)$ and shows that for even $m$ no
Kempe-from-canonical proof route can produce a Hamilton decomposition. Second,
the final odd-modulus D5 gluing proof is formally parity-blind: once an
even-modulus boundary package supplies the same first-exit theorem,
componentwise bridge, regular continuation, and exceptional interface, the odd
proof carries over unchanged. We then isolate a D5 analogue of the Route~E
strategy: the real new work should be a critical-row theorem on the regular
source windows, not a reworking of the final globalization step.
\end{abstract}

\section{What the d=3 even case suggests for D5}

The d=3 even proof has two logically separate pieces. First, a parity barrier
shows that the canonical Kempe route cannot work. Second, a direct low-layer
construction (Route~E) rebuilds the clock-and-carry mechanism and then proves
Hamiltonicity by a finite-defect first-return analysis.

The present D5 odd-modulus theorem already supplies the second half of that
picture at the boundary-theorem level: once the accepted package reaches the
componentwise boundary coordinate
\[
\delta = q + m\sigma,
\qquad
\sigma \equiv w+u-q-1 \pmod m,
\]
one has the exact splice law
\[
\delta' = \delta + 1 \pmod{m^2}
\]
on each splice-connected accessible component, and the remaining globalization
issue collapses to one exceptional row. So for even $m$ the first question is
whether there is a parity obstruction analogous to the d=3 one. The answer is
yes.

\section{A Kempe parity barrier in dimension 5}

We state the sign-product invariant in a form that is independent of the number
of colors.

\begin{definition}[Coloring and Kempe swap]
Let $V$ be a finite set and let $(f_0,\dots,f_{d-1})$ be permutations of $V$.
Fix two distinct colors $r\neq s$ and define the Kempe map
\[
\tau_{r,s}:=f_s^{-1}\circ f_r.
\]
If $X\subseteq V$ is a union of cycles of $\tau_{r,s}$, the \emph{Kempe swap}
of colors $(r,s)$ on $X$ is the new $d$-tuple $(f'_0,\dots,f'_{d-1})$ defined
by exchanging $f_r$ and $f_s$ on $X$ and leaving all other colors unchanged.
\end{definition}

\begin{proposition}[General sign-product invariant under Kempe swaps]
\label{prop:signinv}
Let $(f_0,\dots,f_{d-1})$ be a coloring in the above sense, and let
$(f'_0,\dots,f'_{d-1})$ be obtained from it by one Kempe swap. Then
\[
\prod_{c=0}^{d-1} \sgn(f'_c)
=
\prod_{c=0}^{d-1} \sgn(f_c).
\]
\end{proposition}

\begin{proof}
Only the two swapped colors matter, so it is enough to consider colors $r$ and
$s$. Because the swap support is a disjoint union of $\tau_{r,s}$-cycles and
sign is multiplicative, it is enough to treat one cycle
$C=(v_0,v_1,\dots,v_{\ell-1})$ of $\tau_{r,s}$. Let $\sigma$ be the $\ell$-cycle
$(v_0\,v_1\,\cdots\,v_{\ell-1})$ on $C$ and the identity outside $C$.
By the defining relation $f_r(v_t)=f_s(v_{t+1})$ along $C$, one has
\[
 f'_r = f_r\circ\sigma^{-1},
 \qquad
 f'_s = f_s\circ\sigma.
\]
Therefore
\[
\sgn(f'_r)\sgn(f'_s)
=
\sgn(f_r)\sgn(\sigma^{-1})\,\sgn(f_s)\sgn(\sigma)
=
\sgn(f_r)\sgn(f_s).
\]
All other colors are unchanged, so the full sign product is invariant.
\end{proof}

\begin{corollary}[Parity barrier for Kempe-from-canonical in $D_5(m)$]
\label{cor:d5barrier}
Let $m$ be even. Then no sequence of Kempe swaps starting from the canonical
coloring of the directed torus
\[
D_5(m)=\operatorname{Cay}((\mathbb Z_m)^5,\{e_0,e_1,e_2,e_3,e_4\})
\]
can produce a decomposition into five directed Hamilton cycles.
\end{corollary}

\begin{proof}
Let $N=m^5$. A directed Hamilton cycle on $N$ vertices is an $N$-cycle
permutation, hence has sign $(-1)^{N-1}=-1$ because $N$ is even when $m$ is
even. So any five-color Hamilton decomposition has total sign product
\[
(-1)^5=-1.
\]

In the canonical coloring, each color is translation by one basis vector. Its
cycle decomposition therefore consists of $m^4$ disjoint directed $m$-cycles.
Each $m$-cycle has sign $(-1)^{m-1}=-1$, and since $m^4$ is even one gets
\[
\sgn(f_c)=(-1)^{m^4}=+1
\qquad (0\le c\le 4).
\]
Hence the canonical sign product is $+1$. By
Proposition~\ref{prop:signinv}, this sign product is invariant under all Kempe
swaps, so a Kempe-from-canonical route to a five-cycle Hamilton decomposition
is impossible.
\end{proof}

\begin{remark}[Consequence]
The even-$m$ D5 problem, if true, should be approached as a direct-design or
Route~E-type repair theorem, not as a Kempe-from-canonical continuation of the
odd case.
\end{remark}

\section{The final odd D5 theorem is parity-blind}

The odd-modulus cleanup theorem uses only the following inputs on the true
accessible boundary union:
\begin{enumerate}[label=\textup{(\roman*)},leftmargin=2.5em]
    \item the first-exit theorem with universal targets
    \[
    T_{\mathrm{reg}}=(m-1,m-2,1),
    \qquad
    T_{\mathrm{exc}}=(m-2,m-1,1);
    \]
    \item the componentwise concrete bridge
    \[
    \delta=q+m\sigma,
    \qquad
    \delta' = \delta+1 \pmod{m^2};
    \]
    \item the tail-length reduction at fixed realized $\delta$;
    \item regular continuation for every realized regular label;
    \item the chart/interface landing
    \[
    3m-3 \to 3m-2 \to 3m-1.
    \]
\end{enumerate}
No parity calculation occurs in the final gluing argument itself.

\begin{theorem}[Formal extension theorem]
\label{thm:formalextension}
Let $m\ge 5$ be arbitrary. Assume that on the true accessible boundary union of
an even- or odd-modulus D5 package the following statements hold:
\begin{enumerate}[label=\textup{(\alph*)},leftmargin=2.5em]
    \item \textbf{First exits.} The actual active branch agrees with the
    candidate branch up to first exit, and the regular and exceptional families
    first exit at
    \[
    T_{\mathrm{reg}}=(m-1,m-2,1),
    \qquad
    T_{\mathrm{exc}}=(m-2,m-1,1),
    \]
    with the exceptional first exit leaving in direction $1$.
    \item \textbf{Componentwise bridge.} On each splice-connected accessible
    component, the boundary coordinate
    \[
    \sigma \equiv w+u-q-1 \pmod m,
    \qquad
    \delta := q+m\sigma
    \]
    satisfies
    \[
    \delta' = \delta + 1 \pmod{m^2},
    \]
    and the current boundary event is a function of $(\beta,\delta)$.
    \item \textbf{Tail reduction.} At fixed realized $\delta$, any two
    realizations can differ only by remaining full-chain tail length.
    \item \textbf{Regular continuation.} Every regular realization of every
    realized $\delta$ has a forward successor on the true accessible union.
    \item \textbf{Exceptional interface.} At chart/interface level the
    exceptional continuation is pinned to
    \[
    3m-3 \to 3m-2 \to 3m-1.
    \]
\end{enumerate}
Then every actual lift of the exceptional cutoff row $\delta=3m-3$ continues
through
\[
3m-2 \to 3m-1,
\]
fixed realized $\delta$ determines endpoint class and padded future current-event
word, and raw global $(\beta,\delta)$ is exact on the true accessible boundary
union.
\end{theorem}

\begin{proof}
The proof is exactly the same as in the odd-modulus cleanup theorem, so we only
record the logical chain.

By regular continuation, every regular chain has a forward successor. By the
componentwise bridge, every interior exceptional chain also has a forward
successor because its realized boundary labels follow the odometer step
$\delta\mapsto\delta+1$. Hence only the exceptional cutoff row
$\delta=3m-3$ remains.

For that row, the first-exit theorem identifies the actual exceptional first
exit with the universal target $T_{\mathrm{exc}}=(m-2,m-1,1)$ and fixes the exit
direction. The chart/interface hypothesis then forces the continuation in chain
labels to be
\[
3m-3 \to 3m-2 \to 3m-1.
\]
So the exceptional cutoff also has a forward successor.

Thus every accessible full chain has a forward successor. If some
splice-connected accessible component were finite, its realized boundary image
would be a finite forward orbit segment of the odometer and would therefore
have a terminal right endpoint, contradicting the previous paragraph. Hence
every accessible component is total.

Totality implies that every realized boundary state has infinite remaining tail
length. The tail-reduction hypothesis then shows that fixed realized $\delta$
determines tail length, endpoint class, and padded future current-event word.
Together with the event readout from the componentwise bridge, this is exactly
raw global $(\beta,\delta)$ exactness.
\end{proof}

\begin{corollary}[What the even case really has to prove]
\label{cor:realwork}
To prove the even-$m$ D5 boundary theorem, one does \emph{not} need a new final
gluing argument. One needs an even-modulus package that supplies the same
first-exit law, the same componentwise bridge, regular continuation, and the
same exceptional interface.
\end{corollary}

\section{A D5 analogue of Route~E: critical rows rather than critical lanes}

The bundled D5 interval summaries suggest that the regular source windows obey a
very rigid arithmetic law. We record that law as a programmatic target rather
than as a proved theorem, because the bundled data presently verify it only on
the checked odd-modulus range.

\begin{remark}[Empirical regular-window law from the bundled summaries]
For each regular source parameter $u$, the bundled summaries exhibit a
$\Theta=2$ start label
\[
 s_u = m(u+2)-1 \pmod{m^2},
\]
a regular endpoint label
\[
 e_u = m(u-1)-2 \pmod{m^2},
\]
and a regular-source window of length $m(m-3)$ running from $s_u$ to $e_u$.
The successor of that window is
\[
 e_u+1 = m(u-1)-1 = s_{u-3}.
\]
So the generic source-to-source splice is the cyclic shift
\[
 u \longmapsto u-3 \pmod m.
\]
In the checked odd-modulus data, the unique non-generic source is $u=3$; its
endpoint is the exceptional cutoff row $3m-3$, and the distinguished regular
interface is the source-$4$ terminal / source-$1$ start splice
\[
 3m-2 \to 3m-1.
\]
\end{remark}

This suggests the D5 even proof should imitate d=3 Route~E at the level of
\emph{rows} instead of lanes.

\begin{quote}
Preserve the bulk odometer law $\delta\mapsto\delta+1$, identify the regular
source windows explicitly, prove that their endpoint splice is the generic
shift $u\mapsto u-3$, and isolate the source-$3$ window as the unique finite
splice defect.
\end{quote}

That is the D5 analogue of the d=3 statement ``the bulk motion is a primitive
clock and only finitely many defect lanes remain''.

\begin{proposition}[Critical-row reduction principle]
Assume, for some modulus $m$, that the following hold on the true accessible
boundary union.
\begin{enumerate}[label=\textup{(\roman*)},leftmargin=2.5em]
    \item There is a partition of the regular sector into source windows
    $I_u$, indexed by the regular sources, such that inside each $I_u$ the
    successor law is the generic odometer step $\delta\mapsto\delta+1$.
    \item Each regular window $I_u$ starts at $s_u$ and ends at $e_u$ with
    \[
    s_u = m(u+2)-1,
    \qquad
    e_u = m(u-1)-2,
    \qquad
    e_u+1=s_{u-3}.
    \]
    \item The unique missing source is $u=3$, and its window is the
    exceptional one.
    \item The exceptional window ends at $3m-3$ and its next two labels are
    forced to be $3m-2$ and $3m-1$.
\end{enumerate}
Then the hypotheses of Theorem~\ref{thm:formalextension} are satisfied.
In particular, raw global $(\beta,\delta)$ is exact.
\end{proposition}

\begin{proof}
The generic law on each regular window gives componentwise continuation inside
that window. The identity $e_u+1=s_{u-3}$ gives continuation across all regular
window endpoints, so the regular sector is closed. The source-$3$ hypothesis
isolates the exceptional sector, and the final splice hypothesis gives the
exceptional interface
\[
3m-3 \to 3m-2 \to 3m-1.
\]
Therefore all hypotheses of Theorem~\ref{thm:formalextension} hold.
\end{proof}

\begin{remark}[Interpretation]
This proposition is the closest D5 analogue of the d=3 even ``critical-lane''
idea. It says that the right new theorem for even $m$ is not a fresh global
bridge theorem but a finite-splice theorem for the source windows.
\end{remark}

\section{Bottom line}

Three conclusions seem stable.
\begin{enumerate}[label=\textup{(\arabic*)},leftmargin=2.5em]
    \item The d=3 parity obstruction persists in D5: for even $m$, a
    Kempe-from-canonical route is blocked.
    \item The odd D5 final proof already has the right form for all moduli; it
    becomes an even theorem as soon as the upstream package is rebuilt.
    \item The most promising D5 analogue of Route~E is a \emph{critical-row}
    theorem: preserve the bulk boundary odometer, prove the regular
    source-window splice $u\mapsto u-3$, and isolate source $3$ as the unique
    exceptional splice defect.
\end{enumerate}

So the natural next theorem to attack is not ``reprove the final globalization
step for even $m$'', but rather:
\begin{quote}
\emph{construct an even-modulus D5 boundary package whose only non-generic row
splice is the source-$3$ / exceptional interface.}
\end{quote}

That is the direct D5 analogue of Route~E.

\end{document}
